%% ============================================================================
%% Appendices
%% ============================================================================

\appendices

\section{Stage-to-Citation Map}
\label{app:citations}

Table~\ref{tab:citemap} indexes, for every pipeline stage and
experimental round, the literature it draws on and the role each work
plays --- so that each citation can be audited in place.

\begin{table*}[!t]
\caption{Which Work Is Cited at Which Stage, and Why}
\label{tab:citemap}
\centering
\begin{threeparttable}
\begin{tabular}{@{}p{3.5cm}p{4.1cm}p{9.2cm}@{}}
\toprule
Stage (section) & Works & Role of the citation \\
\midrule
Problem framing (\S\ref{sec:intro}) &
\cite{hilliard2025flourishing,lewis2020rag,gao2023ragsurvey,edge2024graphrag,peng2026graphragsurvey} &
faith-dimension LLM deficits motivating grounded QA; RAG foundations and
surveys; GraphRAG paradigm \\
\addlinespace
Chunking \& granularity (\S\ref{sec:chunking}, \S\ref{sec:audit} M5--M6) &
\cite{chen2024bgem3,edge2024graphrag,lairgi2025atom,liu2024lost} &
tokenizer/embedder fit; chunk size $\to$ extraction density (600 vs.\
2{,}400 tokens $\approx 2\times$ entity references); finer units extract
more completely; long-context degradation as the upper bound \\
\addlinespace
Coreference \& overlap (\S\ref{sec:audit} M4) &
\cite{yao2019docred} &
40.7\% of relational facts cross-sentence, 17.6\% need coreference ---
the structural ceiling of fragment-local extraction without coreference
(used as \emph{indirect} evidence; no peer-reviewed overlap-size
quantification exists) \\
\addlinespace
Entity extraction (\S\ref{sec:entity}) &
\cite{ckip,gemmateam2025gemma3,meher2025corekg,meher2025insidecorekg} &
Chinese segmentation/NER; the local classify-only LLM; structured
prompting as noise suppression \\
\addlinespace
Relation extraction (\S\ref{sec:relation}) &
\cite{bian2025kgsurvey,huang2025tcrqf} &
schema-based paradigm positioning; evidence-span grounding against
context-divorced triples \\
\addlinespace
Cross-references / TSK (\S\ref{sec:crossrefs}) &
\cite{tsk,zheng2025lessismore} &
cross-reference data (OpenBible.info, CC~BY); repair-by-reweighting
rather than deletion \\
\addlinespace
KG-quality remediation roadmap beyond P0, not executed (\S\ref{sec:audit}, \S\ref{sec:future}) &
\cite{meher2025linkkg,zhang2024edc,loconte2026betterlater,zhang2025rakg,zhang2026crossaug,neo4jbuilder,nanographrag} &
coreference recipe ($-45\%$ duplication, length-amplified); post-hoc
canonicalization / correction; cross-chunk reassembly; storage
$\ne$ extraction unit; true-loop gleanings \\
\addlinespace
Retrieval substrate (\S\ref{sec:strategies}) &
\cite{karpukhin2020dpr,robertson2009bm25,cormack2009rrf,chen2024bgem3} &
dense retrieval; BM25 and RRF hybrid (designed, never executed online);
embedder and cross-encoder \\
\addlinespace
Rank fusion (\S\ref{sec:fusionlayer}, \S\ref{sec:ablation}) &
\cite{bacellar2026calibrated} &
calibrated graph/vector fusion as the sophisticated alternative to our
single-$\alpha$ linear form \\
\addlinespace
Evaluation framework (\S\ref{sec:eval}) &
\cite{es2023ragas,zeng2025unbiased} &
RAGAS metrics; judge-bias methodology (position/verbosity;
noncommittal artifact) \\
\addlinespace
Negative result \& benchmark dependence (\S\ref{sec:p0},
\S\ref{sec:round3}, \S\ref{sec:discussion}) &
\cite{zhou2025brink,zeng2025unbiased,xiang2025whengraphs,holm1979} &
KG-defect $\to$ answer non-linearity and defect injection; debiased
gains; when graphs help RAG at all; multiplicity adjustment of the
Round~3 graph $-$ no-graph tests \\
\addlinespace
Graph retrieval on the rebuilt system (\S\ref{sec:rebuiltgraph}) &
\cite{gutierrez2025hipporag2,holm1979,zhou2025brink,xiang2025whengraphs,zeng2025unbiased} &
personalized PageRank as the operator family of the PPR arms (GD,
QD+G, IT-GD; not a reproduction of HippoRAG~2, which ranks passages by
their PPR scores); Holm adjustment where a test family was declared;
degradation under incomplete knowledge (mechanism~(B));
task-dependent graph benefit and debiased GraphRAG gains
(mechanism~(C)) \\
\addlinespace
Aggregation, event-temporal layer \& graph-native selection (proposed,
not built) (\S\ref{sec:diagnosis}, \S\ref{sec:future}) &
\cite{sarthi2024raptor,edge2024graphrag,traag2019leiden,lairgi2025atom,zhang2025entityevent,sun2025dygrag,gutierrez2025hipporag2,wang2025archrag,bacellar2026calibrated} &
RAPTOR trees / Leiden and hierarchical community summaries for
cross-chapter aggregation (untested); temporal--causal event modeling
(the event-temporal layer, no longer planned); passage ranking by PPR
scores as a selector for bridge passages, as in HippoRAG~2, which the
rebuilt-system experiments did not reproduce, or calibrated
graph--vector fusion \\
\addlinespace
Retrieval-side chunking sensitivity (\S\ref{sec:related}) &
\cite{shaukat2026chunking} &
cited \emph{only} for retrieval-downstream chunking sensitivity ---
explicitly not extrapolated to KG construction \\
\bottomrule
\end{tabular}
\end{threeparttable}
\end{table*}

\section{Reproducibility Notes}
\label{app:repro}

\textbf{Determinism and re-scoring.} All retrieval metrics are computed
programmatically against the gold references (at unit, verse, and
chapter level); every
evaluation run archives per-question raw responses (sources,
strategies, raw reranker and fused scores, route), and the metric suite
can be recomputed bit-exactly from the archives. The fast
retrieval-only loop was validated to agree bit-exactly with the full
pipeline on all retrieval metrics in Rounds~0--2, while the intent
classifier was inactive. Retrieval is deterministic \emph{given the
classifier output}; routes and graph inputs depend on the LLM intent
classifier, which samples at temperature 0.1 --- in Round~3, 12 of 500
questions (2.4\%) were routed differently in two configurations whose
routing logic is identical. The script \code{paper/tools/round3\_stats.py} recomputes
every number derived from the three Round~3 archives (aggregates,
bootstrap intervals with a fixed seed, the noise floor, the slot
analysis, the faithfulness gate, the traces) after asserting that its
aggregates equal the archived ones; \code{paper/tools/replay\_routes.py}
reproduces the empty-classifier route replay and the July~15
comparisons from the git-archived runs, but it imports backend modules
that the October 2026 rebuild removed, so it must be run from a worktree
checked out at a legacy commit.

\textbf{Comparability across the three time points.} Answer model
(Gemma~4 E4B, 8-bit), judge model (Gemma~4 26B), decoding temperature
(0.1), and top-$k$ (5) were verified identical, as was the route
distribution (R1$=$20, R2$=$25, R3$=$21, R4$=$18, R5$=$10, R6$=$4,
fallback$=$2) at the first two time points; at the third, one question
(EVENT\_015) moved from R6 to R4 after the event keyword ``Gethsemane
prayer'' was added in the fusion round. All three time points ran with
the intent classifier silently inactive (Section~\ref{sec:signals}):
replaying the signal detector with an empty classifier output
reproduces every logged route. The entity-query strategy was off at the
first two time points and on after the fusion round (its re-enablement
is itself one of the audited repairs). The only other systematic
difference between time points is the state of the knowledge graph and
of the post-retrieval chain, as described in the text.

\textbf{Comparability with Round 3.} Round~3 differs from the three
time points in the classifier (functioning), book-anchor retrieval (per
book, developed and checked on the 500 questions), merged-verse lookup,
the judges' inputs (generator context blocks) and faithfulness rubrics,
the Ollama server version, and the benchmark (500 questions). The bridge is the
original 100 questions scored by the deterministic retrieval metrics,
which can be recomputed from every archive; across it the graph
configuration's hit rate (0.97) and its margin over semantic-only
($+0.16$) are unchanged.

\textbf{Rollbackability.} While the legacy system was deployed, every P0
and fusion-round mutation was reversible: backfilled mention edges,
co-occurrence relation edges, and TSK cross-reference edges carried
provenance flags and could be deleted by a single predicate each;
curated event nodes carried a source tag; alias injections had
before-snapshots; and the rank-fusion layer, expansion caps, and
entity-query strategy were configuration switches whose legacy settings
reproduced the pre-remediation system exactly. Neo4j and the legacy data
were deleted in the October 2026 rebuild (Appendix~\ref{app:rebuild}),
so the legacy graph and service can now be restored only from the
pre-rebuild backup.

\textbf{Rebuilt-system graph experiments.} The new graph experiments of
Section~\ref{sec:rebuiltgraph} (GD, realkg, iterq, the per-hop probes,
and the October 18 run) ran offline on the data layer of Rebuild-2; GN-41 ran on the legacy KG
before the rebuild, and the event-lane and release comparisons on the
served builds. The pre-registrations and code of GD and the October 18
run are in the repository:
\code{evaluation/experiments/2026-10-14\_graph\_gd/} (\code{prereg.md};
\code{evaluate.py}, whose subcommands recompute \code{new82\_report.json},
\code{v1\_old31.json}, and \code{units\_old31.json}; \code{v1\_power.py},
which simulates the stop rule's operating characteristics) and
\code{evaluation/experiments/2026-10-18\_graph\_qd/} (\code{plan.md},
committed in 16f9ced before the run started and updated with the results
in 533b4ee; \code{qd\_eval.py}; \code{qd\_step2.py}). The archived
outputs are kept outside the repository, in the project's artifact
store under \code{reports/}:
\code{graph\_gd\_20261014/}; \code{graph\_qd\_20261018/}
(\code{report.json}, \code{step2/report2.json});
\code{kg\_effect\_eval\_20261014/new82\_ans/} (\code{retrieval\_iter.json},
the iterative arms); \code{multihop\_kg\_20261012/} (the \code{realkg} and
\code{iterq} analyses on OLD31, each with its \code{analyze.py}, and the
per-hop probes); \code{gn41\_20261002/} (\code{stats/analyze.py});
\code{eval3c\_20261014\_analysis/} (\code{analyze\_eval3c\_v2.py},
\code{final\_out/analysis.json}); and \code{r2eval/} (the Rebuild-2
gates). In the October 18 run, 32 reference reranker pairs were
re-scored before the run and after each of its two steps, with a
maximum absolute difference of 0. The Holm adjustment of Round~3
(Section~\ref{sec:round3}) is archived as output only
(\code{verify\_stats/multi.out} in the October 10 evidence snapshot
\code{tmp\_evidence\_20261010/} under the store's \code{reports/}, with
the test function in \code{vcommon.py}); the driver script that produced
it was not archived, while \code{paper/tools/round3\_stats.py} recomputes
the unadjusted intervals.

\textbf{Scale of the evaluation.} A full 100-question pipeline run
takes $\approx 2.5$ hours on the single-host deployment (36 min
collection, 1\,h\,42\,m RAGAS judging, 13 min coverage judging); the
retrieval-only loop takes $\approx 15$ minutes and was used for the
$\alpha$ sweep and regression forensics. On the 500-question benchmark,
LLM judging alone took roughly 10--12 hours per Round~3 configuration
(two faithfulness judges, three further RAGAS metrics, and coverage).

\section{Post-Audit Corrections and the October 2026 Rebuild}
\label{app:rebuild}

An October 2026 audit of the data pipeline contradicted several
statements of earlier versions of this paper. This revision corrects
them in place and leaves every experimental number of Rounds~0--3, which
remain the run of record, unchanged: the corpus is the RCUV, not the
new-punctuation CUV named on its PDFs, and its legacy conversion lost
text and added ghost verses (Section~\ref{sec:preprocessing}), so the
legacy store's 31{,}102 verse positions are not the PDF's 31{,}092
(Section~\ref{sec:crossrefs}); the OpenBible.info TSK data are CC~BY, not
public domain (Section~\ref{sec:crossrefs}); hybrid retrieval never ran,
so every round is dense-only (Section~\ref{sec:strategies}); and the
serving-side tokenization of Rounds~0--3 (the embedder's query encoding
and the reranker's query--passage input) was inconsistent with the build
side, a defect whose correction had no measurable retrieval effect
(Section~\ref{sec:evalinfra}). The audit also measured the legacy
knowledge graph of Section~\ref{sec:kg}: it had no entity-identity
layer --- an \code{entity\_id} was the type plus toneless pinyin, so
homonyms and homophones shared one node (of the 175 edges on the apostle
John's node, 29 concern the apostle) --- and the strict precision of its
15{,}926 semantic edges, called fact edges in earlier versions, was only
about 0.23--0.28 (lenient 0.34--0.45), with most wrong edges in the
co-occurrence backfill (Section~\ref{sec:relation}). The system has
since been rebuilt from the
66 PDFs in versioned layers. The rebuilt text matches the PDFs with zero
character-conservation residual: 31{,}021 verse labels (70 of them merged
verses such as ``29--30'') over 31{,}092 verse positions, 11 verse numbers
that the RCUV gives only in footnotes (the legacy ghost verses among
them), all 116 psalm superscriptions, and 1{,}013 footnotes. Sparse
retrieval is retired and the serving-side tokenization corrected.

Two rebuilt builds have been served: Rebuild-1
(\code{b20261008\_6daa4f31}) and Rebuild-2 (\code{b20261008\_e05d3e55}),
the build now in service. (We avoid the labels R1/R2 for builds, since
R1--R6 name routes.) Neither uses Neo4j. The knowledge layer is reduced
to an L0 name layer (2{,}926 names), which the service uses only as
routing vocabulary for person and place signals, and an L1 curated event registry
(31 events in Rebuild-2; 33 in Rebuild-1, whose registry was exported
from the curated events of the legacy graph); there is no
entity-identity layer and no relation layer. The entity, relation, and
cross-reference graph of Section~\ref{sec:kg} was retired rather than
rebuilt, and TSK is not loaded. The registry is a curated event lookup
lane, not a graph operation: a literal trigger match appends at most one
anchor passage after the completed top-$k$ on the event and cross-book
routes (R4, R5), and the appended passage takes no part in ranking. In a
three-arm run on the 500 GT~v2 questions (Section~\ref{sec:rebuiltresults};
one generation per arm), turning the lane on changed answer coverage by
$+0.0089$ ($[-0.0030, +0.0207]$). The lane fired on only 23 of the 500
questions, and on those it adds a sixth passage, so the comparison is
not $k$-aligned; we report it as descriptive, not as graph evidence.
Rebuild-2 likewise showed no measurable retrieval gain over Rebuild-1
(verse recall@6 $+0.0002$, $[-0.0016, +0.0026]$; descriptive). In the
same three-arm run, the full pipeline beat semantic-only retrieval in
answer coverage by $+0.032$ ($[+0.012, +0.053]$; descriptive), a gain
concentrated in the 179 questions routed to the exact-verse and chapter
routes (R1, R2: $+0.092$, $[+0.048, +0.138]$), that is, in citation
resolution; on the other 321 questions no gain was detected ($-0.002$,
upper bound $+0.017$). The
graph-centric serving architecture of
Sections~\ref{sec:kg}--\ref{sec:retrieval} is therefore historical
(Rounds~0--3); the graph-retrieval experiments on the rebuilt system are
reported in Section~\ref{sec:rebuiltgraph}.
